% Problem record op_4b453163f7f8c748. This TeX file is derived from
% database/problems_json/op_4b453163f7f8c748.json by scripts/sync-tex.mjs; edit the JSON record
% and rerun the script rather than editing this file.
\section{LOCC entanglement cost of PPT-entangled Gaussian states}
\label{sec:op_4b453163f7f8c748}

\paragraph{Problem.}
What is the exact LOCC entanglement cost of a finite-energy PPT-entangled Gaussian state?

Let $\rho_V$ be a zero-mean Gaussian state of $m$ modes held by Alice and $n$ modes held by Bob, with integers $m,n\geq2$. Assume that $\rho_V$ is entangled and its partial transpose $\rho_V^{T_B}$ is positive. Its finite covariance matrix and canonical quadratures satisfy

\begin{equation}
\begin{gathered}
V_{jk}:=\operatorname{Tr}\rho_V\{R_j,R_k\},\qquad
[R_j,R_k]=i(\Omega_{m+n})_{jk},\\
\Omega_k:=\bigoplus_{j=1}^k\begin{pmatrix}0&1\\-1&0\end{pmatrix}.
\end{gathered}
\label{eq:ser15-statement-1}
\end{equation}

Equation~\eqref{eq:ser15-statement-1} uses vacuum covariance $I$. The allowed channels $\Lambda_N$ use unrestricted local operations and classical communication (LOCC). Define the Bell-pair density operator by

\begin{equation}
\Phi_2:=|\phi_2\rangle\langle\phi_2|,\qquad
|\phi_2\rangle:=(|00\rangle+|11\rangle)/\sqrt2.
\label{eq:ser15-statement-2}
\end{equation}

Using Eq.~\eqref{eq:ser15-statement-2}, the entanglement cost $E_C(\rho_V)$ is the infimum of rates $r\geq0$ for which LOCC channels $\Lambda_N$ satisfy

\begin{equation}
\lim_{N\to\infty}\|\Lambda_N(\Phi_2^{\otimes\lceil rN\rceil})-\rho_V^{\otimes N}\|_1=0.
\label{eq:ser15-statement-3}
\end{equation}

Determine the cost specified by Eq.~\eqref{eq:ser15-statement-3} as a function of $V$. Separability means the trace-norm closed convex hull of product states. Here $\|X\|_1:=\operatorname{Tr}\sqrt{X^\dagger X}$.

\subsection*{Status}

Unsolved

\subsection*{Source}

This exact-evaluation question combines the PPT-entangled Gaussian family of Werner--Wolf, Section IV, with the infinite-dimensional operational cost theorem \sourcecite{ref:ser15-1}{WW01}, \sourcecite{ref:ser15-2}{YKHL25}.

\subsection*{Progress}

\begin{itemize}
  \item Positive partial transpose does not imply Gaussian separability once both parties have two modes, and the logarithmic negativity vanishes on every state in this question:

\begin{equation}
\rho_V^{T_B}\geq0\quad\Longrightarrow\quad
\log_2\|\rho_V^{T_B}\|_1=0.
\label{eq:ser15-progress-1-1}
\end{equation}

Despite Eq.~\eqref{eq:ser15-progress-1-1}, explicit four-mode counterexamples to the converse separability implication exist. \sourcecite{ref:ser15-1}{WW01}

  \item Theorem 7 of Yamasaki et al. applies because the local entropies are finite. Finite-energy Gaussian states satisfy the infinite-dimensional entanglement-cost formula

\begin{equation}
\begin{gathered}
E_C(\rho_V)=\lim_{N\to\infty}\frac{E_F(\rho_V^{\otimes N})}{N},\\
E_F(\omega):=\inf_{\mu:\,\int|\psi\rangle\langle\psi|\,d\mu(\psi)=\omega}
\int S(\operatorname{Tr}_B|\psi\rangle\langle\psi|)\,d\mu(\psi),
\end{gathered}
\label{eq:ser15-progress-2-1}
\end{equation}

In Eq.~\eqref{eq:ser15-progress-2-1}, $\mu$ ranges over probability measures on normalized pure states and $S(\tau):=-\operatorname{Tr}(\tau\log_2\tau)$; this variational formula permits non-Gaussian decompositions. \sourcecite{ref:ser15-2}{YKHL25}
\end{itemize}

\subsection*{References}

\begin{enumerate}[label={},leftmargin=4.8em,labelsep=0.5em]
  \footnotesize
  \item[\textup{[WW01]}]\label{ref:ser15-1}
  R. F. Werner and M. M. Wolf, "Bound Entangled Gaussian States," \emph{Physical Review Letters} \textbf{86}, 3658--3661 (2001). \href{https://doi.org/10.1103/PhysRevLett.86.3658}{doi:10.1103/PhysRevLett.86.3658}; \href{https://arxiv.org/abs/quant-ph/0009118}{arXiv:quant-ph/0009118}.

  \item[\textup{[YKHL25]}]\label{ref:ser15-2}
  H. Yamasaki, K. Kuroiwa, P. Hayden, and L. Lami, "Entanglement Cost for Infinite-Dimensional Physical Systems," \emph{Communications in Mathematical Physics} \textbf{406}, 277 (2025). \href{https://doi.org/10.1007/s00220-025-05431-1}{doi:10.1007/s00220-025-05431-1}; \href{https://arxiv.org/abs/2401.09554}{arXiv:2401.09554}.
\end{enumerate}

\subsection*{Comment}

The regularized convex-roof characterization is established. Its exact evaluation for arbitrary PPT-entangled Gaussian covariances remains open. Replacing the unrestricted roof by a Gaussian roof or changing LOCC to PPT-preserving operations changes the question.

\subsection*{Field}

Quantum Resource Theory

\subsection*{Topic}

Gaussian quantum information; Bound entanglement; Entanglement cost

\subsection*{ID}

\texttt{op\_4b453163f7f8c748}
